\documentclass[lettersize,journal]{IEEEtran}
\usepackage{amsmath,amsfonts}
\usepackage{algorithmic}
\usepackage{algorithm}
\usepackage{array}
\usepackage[caption=false,font=normalsize,labelfont=sf,textfont=sf]{subfig}
\usepackage{textcomp}
\usepackage{stfloats}
\usepackage{url}
\usepackage{verbatim}
\usepackage{graphicx}
\usepackage{cite}
\hyphenation{op-tical net-works semi-conduc-tor IEEE-Xplore}
% updated with editorial comments 8/9/2021

\begin{document}

\title{An IoT System for Real-time Human Activity Recognition}

\author{Tianchen Liu$^{*}$, Juntao Hu$^{*}$
        % <-this % stops a space
% \thanks{This paper was produced by the IEEE Publication Technology Group. They are in Piscataway, NJ.}% <-this % stops a space
% \thanks{Manuscript received April 19, 2021; revised August 16, 2021.}
\thanks{$^*$ equal contribution.}}
% The paper headers
% \markboth{Journal of \LaTeX\ Class Files,~Vol.~14, No.~8, August~2021}%
% {Shell \MakeLowercase{\textit{et al.}}: A Sample Article Using IEEEtran.cls for IEEE Journals}

% \IEEEpubid{0000--0000/00\$00.00~\copyright~2021 IEEE}
% Remember, if you use this you must call \IEEEpubidadjcol in the second
% column for its text to clear the IEEEpubid mark.

\maketitle

\begin{abstract}

This project builds an IoT system for real-time Human Activity Recognition (HAR) that describes human activities with natural language captions. A lightweight model runs locally on the smartphone and immediately produces an activity caption. Selected frames are then uploaded to a server, where a stronger
Vision-Language Model (VLM) verifies the local caption and either accepts or rewrites it. However, VLM verification must share an end-to-end latency budget between wireless transmission and server computation, while the VLM execution time varies with the input and the current compute load. We therefore introduce a lightweight latency predictor that estimates the
edge-side verification time from the image size, current CPU utilization, and
recent VLM execution latency. The remaining end-to-end deadline is converted
into a communication budget and used by an application-layer scheduler to
prioritize verification uploads from multiple clients.

\end{abstract}


\section{Motivation}


An interactive or embodied system has to respond while the action is still
in progress. A closed label such as ``sitting'' is too coarse for that
response. Sitting still, sitting and looking at a phone, and sitting and
waving call for different responses, yet they collapse to one class if the
recognizer only has a small label set. We therefore describe the activity
with an English sentence. Running a strong vision-language model on every captured frame is impractical
for an interactive IoT system because the model is computationally expensive.
Our design therefore utilizes a small model on the smartphone which can write that sentence immediately, and a server with a stronger Vision-Language Model (VLM) can check the sentence and produce a more accurate description. However, sending every picture to the server makes the user wait on the uplink. The system when the picture has to be uploaded to the server. To achieve this, only a
low-confidence frame is JPEG-encoded and sent to the server, where the confidence of
that sentence is the minimum probability among the generated tokens. The server then
checks the picture and replaces the sentence when the two descriptions disagree.

Besides, checking a selected frame has an end-to-end confirmation deadline. That
deadline is shared by wireless transmission and server computation. A fixed
communication budget therefore cannot reliably guarantee an end-to-end
verification deadline. We therefore investigate a compute-aware communication scheduler. Before a
verification frame is uploaded, the system predicts its VLM latency
and subtracts the predicted computation time from the end-to-end deadline.
The remaining time becomes the communication budget of that frame. When
multiple clients request verification simultaneously, these dynamically
derived communication deadlines are used to determine their upload order.
\section{Literature Study}

\subsection{Human Activity Recognition}
Activity recognizers take RGB, skeleton, IMU, radar, or Wi-Fi, and usually
return one label from a closed set. An open-set model can reject an unknown
action. It does not describe that action. An open-vocabulary model can match
a clip to a class name while the action is still going on. Free-form
captions also exist. OV-HAR turns a sensor window into a sentence
\cite{ray2025ovhar}, and OpenHOUSE calls a vision-language model for a
free-form description of streaming video \cite{kang2025openhouse}. Both are
evaluated on recorded data, not on a live phone link. The missing piece is a
sentence produced on the phone, before any server round trip, while the
person is still moving.
\subsection{Vision-Language Models}
CLIP retrieves a sentence from a fixed bank by cosine similarity. It does
not write a new sentence \cite{radford2021learning}. A generative
vision-language model writes the sentence token by token, and accuracy
follows model size. Qwen2-VL scales this family up to 72B parameters
\cite{wang2024qwen2vl}. Models that can sit on a phone also exist.
MiniCPM-V is in the 2B--8B range \cite{yao2024minicpm}, and SmolVLM-256M
fits in under 1~GB \cite{marafioti2025smolvlm}. Published scores for these
models cover documents, charts, and general questions, not fine differences
between similar activities. A small model can write the sentence. It does
not, by itself, say which of those sentences are safe to leave on the
screen.
\subsection{Client--Server Collaboration}
Neurosurgeon profiles network layers and chooses where to split them, so
the server receives image features and the user waits for that transfer
\cite{kang2017neurosurgeon}. Speculative decoding splits the work by
tokens: a small model drafts, a large model verifies, and the server still
reads the picture \cite{leviathan2023fast}. edgeVLM answers on the device
at once and reuses a late cloud caption as context for later frames, rather
than as the result of the frame that is on screen now
\cite{qian2025edgevlm}. EdgeBoost uploads an input only when the local
confidence is low. On image classification this removes about 20--55\% of
the communication of a cloud-only run, and the cloud still returns a class
label \cite{said2024edgeboost}. What these systems do not do is show an
activity sentence immediately and send the picture only when that sentence
is unsafe.
\subsection{Communication}
A proportional-fair scheduler and an L4S queue assign service to packets
that have already been formed. They do not look at how long the laptop will
spend on the picture after it arrives
\cite{jalali2000data,briscoe2023l4s}. Semantic communication scores a
transmission by task success rather than by throughput, but the payload is
still a feature tensor \cite{xie2021deep}. Offloading chooses local or
remote execution, and adaptive streaming changes how many bytes are sent.
After the sentence has already been written, these mechanisms still do not
decide whether the picture should follow, how large that picture should be,
and which phone should be checked first.
Our system takes that last decision. Token confidence decides whether the
JPEG is uploaded. JPEG quality and resolution decide its size. For every
upload, the predicted edge latency sets the uplink deadline, and the
application-layer scheduler serves the earliest deadline first.

\section{Technical Approach}

\subsection{Pipeline}

Let $\mathcal{C}$ be a bank of English activity sentences, each with a CLIP
text embedding $e_c$. For a captured frame $x_i$, SmolVLM-256M runs on the
smartphone in 4-bit form and writes a short action phrase $s_i$ together
with the probability $p_{i,t}$ of each generated token. The confidence of
the phrase is
\begin{equation}
\kappa_i
=
\min_t p_{i,t}.
\end{equation}
The phrase is shown as soon as this local inference finishes. It does not
wait for the server.

The JPEG is uploaded only when the phrase is unsafe:
\begin{equation}
\kappa_i < \tau_{\mathrm{u}}.
\end{equation}
Otherwise the phone sends the sentence alone. CLIP on the server maps that
sentence into the bank,
\begin{equation}
c_{d,i}
=
\arg\max_{c\in\mathcal{C}}
\cos\big(f_T(s_i),e_c\big),
\end{equation}
and this match is the caption that stays on screen.

When the JPEG is uploaded, CLIP also retrieves a caption from the picture,
\begin{equation}
c_i^\star
=
\arg\max_{c\in\mathcal{C}}
\cos\big(f_I(x_i),e_c\big).
\end{equation}
The phone sentence is kept when $c_i^\star=c_{d,i}$. Otherwise the displayed
caption is replaced by $c_i^\star$. A larger generative model, Qwen2-VL, can
take the same place as CLIP: it writes its own phrase, the phrase is mapped
into $\mathcal{C}$ by the same text encoder, and the caption is replaced
when that match differs from $c_{d,i}$.

Only the low-confidence path puts an image on the air:
\[
\mathrm{Smartphone}
\rightarrow
\mathrm{Wi\mbox{-}Fi}
\rightarrow
\mathrm{server~checker}
\rightarrow
\mathrm{Smartphone}.
\]

\begin{algorithm}[t]
\caption{Local sentence and edge check}
\label{alg:verify}
\begin{algorithmic}[1]
\STATE Run 4-bit SmolVLM on frame $x_i$ and display $s_i$.
\STATE $\kappa_i\leftarrow\min_t p_{i,t}$.
\IF{$\kappa_i\geq\tau_{\mathrm{u}}$}
    \STATE Send $s_i$ only and keep its bank match $c_{d,i}$.
\ELSE
    \STATE Predict the edge checking time $\hat T_i^{\mathrm{edge}}$.
    \STATE Compute the uplink budget $B_i^{\mathrm{up}}$.
    \STATE Wait for an upload grant, then send the JPEG.
    \STATE Retrieve $c_i^\star$ with CLIP, or with Qwen2-VL.
    \IF{$c_i^\star=c_{d,i}$}
        \STATE Accept $c_{d,i}$.
    \ELSE
        \STATE Replace the caption with $c_i^\star$.
    \ENDIF
\ENDIF
\end{algorithmic}
\end{algorithm}

\subsection{Latency Model}

The sentence that appears first is local:
\begin{equation}
T_i^{\mathrm{draft}}
=
T_i^{\mathrm{local}},
\end{equation}
where $T_i^{\mathrm{local}}$ is the SmolVLM execution time on the phone.

A low-confidence frame also waits for a checked caption. Its confirmation
delay is
\begin{equation}
T_i^{\mathrm{conf}}
=
T_i^{\mathrm{wait}}
+
T_i^{\mathrm{up}}
+
T_i^{\mathrm{edge}}
+
T_i^{\mathrm{down}},
\end{equation}
where $T_i^{\mathrm{wait}}$ is the time until the scheduler grants the
upload, $T_i^{\mathrm{up}}$ is the JPEG upload time,
$T_i^{\mathrm{edge}}$ is the server checking time, and
$T_i^{\mathrm{down}}$ is the return time of the result. The downlink carries
an acceptance decision or a caption identifier, so the large terms are the
uplink and the checker. The request meets its deadline when
\begin{equation}
T_i^{\mathrm{conf}}
\leq
D_i^{\mathrm{conf}}.
\end{equation}

\subsection{Online Edge-Latency Prediction}

Checker time is not treated as a constant. It changes with the uploaded
image and with the load on the server. The predictor estimates that time
before the JPEG is uploaded, from
\begin{equation}
\mathbf{z}_i
=
\left[
L_i,\;
U_i^{\mathrm{CPU}},\;
T_{i-1}^{\mathrm{edge}}
\right],
\end{equation}
where $L_i$ is the JPEG size in bytes, $U_i^{\mathrm{CPU}}$ is the server
CPU utilization when the request arrives, and $T_{i-1}^{\mathrm{edge}}$ is
the measured checking time of the previous request. The JPEG size is known
once the phone has encoded the frame. It sets the upload time and is also a
cue to image complexity. CPU utilization is read on the server at the
request. The previous checking time carries warm-up, contention, and
model-dependent variation.

A small network $g_\theta(\cdot)$ predicts
\begin{equation}
\hat T_i^{\mathrm{edge}}
=
g_\theta(\mathbf{z}_i).
\end{equation}
It is trained on traces from the server used in the demo. The network stays
small because it supports scheduling, not captioning. Accuracy is the mean
absolute error between predicted and measured checking time.

\subsection{Compute-Aware Link Control}

Three controls sit above the Wi-Fi MAC. The upload threshold
$\tau_{\mathrm{u}}$ decides whether a picture follows the sentence. JPEG
quality and resolution decide how large that picture is. The scheduler
decides which phone is checked first.

For an upload request $i$, the remaining uplink budget is
\begin{equation}
B_i^{\mathrm{up}}
=
\left[
D_i^{\mathrm{conf}}
-
\hat T_i^{\mathrm{edge}}
-
\hat T_i^{\mathrm{down}}
-
M
\right]^+,
\end{equation}
where $\hat T_i^{\mathrm{down}}$ is the expected return time and $M$ is a
margin for prediction error. If the request reaches the scheduler at time
$a_i$, its uplink completion deadline is
\begin{equation}
d_i^{\mathrm{up}}
=
a_i+B_i^{\mathrm{up}}.
\end{equation}
The scheduler serves the pending request with the earliest deadline:
\begin{equation}
i^\star
=
\arg\min_{i\in\mathcal{Q}}
d_i^{\mathrm{up}}.
\end{equation}
Each phone first sends a short request that contains the JPEG size and the
confirmation deadline, then waits for an application-layer grant. The
granted phone uploads the JPEG. A long predicted edge time leaves less room
for the upload, so that request gets an earlier uplink deadline. A short
predicted edge time can wait longer without missing
$D_i^{\mathrm{conf}}$.
\section{Implementation}

\subsection{Platform}

The client is a smartphone browser. SmolVLM-256M-Instruct runs there in
4-bit form: the vision encoder and the decoder use 4-bit weights, on WASM
on an iPhone and on WebGPU where the browser provides it. The generated
sentence and its minimum token probability are available before any image
leaves the phone.

Frames that fall below $\tau_{\mathrm{u}}$ are JPEG-encoded and sent over
Wi-Fi to a laptop. The laptop server is FastAPI. It holds the activity
sentence bank and runs CLIP as the checker. Qwen2-VL is the stronger
generative checker in the real-world comparison. Before the JPEG is
transmitted, the phone sends the planned JPEG size and the confirmation
deadline. The laptop reads its CPU utilization and the previous checking
time, predicts $\hat T_i^{\mathrm{edge}}$, and places the request in the
upload scheduler. The scheduler returns a grant. The phone then uploads the
JPEG, the laptop checks it, and the accept-or-rewrite result comes back.
This path does not change the Wi-Fi MAC or the access point.

\subsection{What is Already Running}

The current prototype already runs this client--server path in the browser.
SmolVLM writes the action phrase on the phone, the page shows that phrase
and the minimum token probability, and the JPEG is uploaded only when the
probability is below the threshold. CLIP on the laptop maps both the phone
sentence and, when present, the picture into the sentence bank. The caption
is replaced only when those two matches differ. A confident phrase sends
the sentence and leaves the picture on the phone.

The predictor, the dynamic uplink budget, and the multi-client scheduler
are not in this prototype yet. The next step is to timestamp the Wi-Fi
transfer and to record checker latency under different CPU loads, JPEG
sizes, and checker models. Those traces train the predictor and drive the
scheduler simulation.

\subsection{Evaluation}

The caption path is tested on the phone and the laptop over the same Wi-Fi
network. One person performs the activities in the sentence bank. For this
real-world test we report on-device generation latency, the fraction of
frames that upload a JPEG, the fraction of captions that are accepted or
rewritten, and a human check of the final caption on a small labeled
subset. The three caption baselines are the phone sentence alone, an upload
on every frame, and the confidence-gated upload.

The scheduler is evaluated in simulation driven by the measured upload and
checking traces, with several phones contending for the laptop. We report
prediction mean absolute error, the fraction of checks that meet the
end-to-end deadline, the 95th percentile confirmation delay, the upload
waiting time, and the same figures as the number of clients grows. The
three schedulers are:
\begin{enumerate}
\item \emph{Round robin}: uploads are served with no deadline and no
checking time;
\item \emph{Static budget}: uplink deadlines use one fixed average checking
time;
\item \emph{Predictive budget}: uplink deadlines use
$\hat T_i^{\mathrm{edge}}$.
\end{enumerate}
The caption comparison and the scheduler comparison are reported
separately.

\subsection{Roles}

Both members integrate the demo, the evaluation, and the report.

Juntao Hu is responsible for the local model, the caption bank, the
acceptance rule, the laptop checker, and the caption-quality evaluation.

Tianchen Liu is responsible for the latency predictor, the dynamic uplink
budget, the upload scheduler, and the multi-client simulation.

The sentence-and-check path is the information-processing part of the
project. Predictive, deadline-aware upload is the networking part.
\section{Project Timelines}

The smartphone SmolVLM path and the laptop CLIP check are the starting
point. Latency prediction, the upload scheduler, and the evaluation will be
finished before the final presentation on 20 November 2026.

\begin{table}[t]
\caption{Schedule}
\label{tab:schedule}
\centering
\begin{tabular}{ll}
\hline
Dates & Work \\
\hline
23 Sep.--3 Oct. &
Local inference and the \\
&
Wi-Fi checking path \\

4--17 Oct. &
Checker latency traces and \\
&
predictor training \\

18--31 Oct. &
Dynamic budget and the \\
&
application-layer upload scheduler \\

1--12 Nov. &
Multi-client simulation and \\
&
scheduler baselines \\

13--20 Nov. &
Final evaluation, demo, report, \\
&
and presentation \\
\hline
\end{tabular}
\end{table}

\bibliographystyle{IEEEtran}
\bibliography{main}

\end{document}
